\documentclass[11pt]{article}
\usepackage{fontspec,amsmath,amssymb,amsthm,hyperref}
\usepackage[margin=1in]{geometry}
\setlength{\emergencystretch}{3em}
\title{Exact Gaussian fourth moments for the Hessian correction}
\author{LangevinLean mathematical exposition}
\date{}
\begin{document}\maketitle
\noindent Source module: \texttt{KineticLangevin/GaussianFourth.lean}. Compile this proof document with XeLaTeX or LuaLaTeX. The accompanying HTML pairs the mathematical argument with the exact Lean statements and records the scope of the formalization.


The Hessian-Lipschitz Taylor error is quadratic in the position increment. Its L\ensuremath{{}^{2}} bound therefore needs a fourth moment. The exact standard Gaussian value is \(p^2+2p\); retaining both terms is necessary in dimension one. These identities are proved for the actual Gaussian measures and transferred to any realization by its proved law.


\section{Differentiate the actual scalar moment-generating function four times}


A centered Gaussian of variance \(v\) has moment-generating function \(e^{vt^2/2}\). Successive derivatives are \(vt e^{vt^2/2}\), \((v+v^2t^2)e^{vt^2/2}\), and \((3v^2t+v^3t^3)e^{vt^2/2}\). The fourth derivative at zero is \(3v^2\). Lean constructs each derivative from the power, multiplication and exponential rules, then invokes the proved theorem identifying the derivative of a Gaussian moment-generating function with its moment. This also includes variance zero.



\noindent Lean declarations: \texttt{centeredGaussian\_fourth\_moment}.

\section{Separate diagonal fourth moments from independent cross moments}


For \(G=(G_1,\ldots,G_p)\), expand \(\lVert G\rVert^4=\sum_{i,j}G_i^2G_j^2\). A diagonal term has expectation three by the scalar theorem. Distinct coordinates are independent under the constructed product Gaussian measure, so their squared product has expectation \(1\cdot1=1\). Finite integrability follows from fourth moments and \(G_i^2G_j^2\le G_i^4+G_j^4\). The sum is \(3p+p(p-1)=p^2+2p\). Lean first works on the actual product of scalar Gaussians, then uses its proved Euclidean standard-Gaussian pushforward identity. Empty coordinate sets are included.



\noindent Lean declarations: \texttt{stdGaussian\_norm\_fourth\_moment}.

\section{Transfer the moment to the coupled stationary velocity}


If an actual random vector has the standard Gaussian law, the pushforward integration identity gives the same norm fourth moment on its original probability space. Applying this to each stationary kinetic velocity supplies exactly the \(\sqrt{p^2+2p}\) factor in the second-order Taylor remainder. No independence between this velocity and the numerical position discrepancy is assumed.



\noindent Lean declarations: \texttt{gaussian\_vector\_fourth\_moment}.

\end{document}
